% concatenated from erdos670.tex
\documentclass[twoside, reqno, 12pt]{amsart}

\usepackage[T1]{fontenc}
\usepackage[utf8]{inputenc}
\usepackage{lmodern}
\usepackage{microtype}
\usepackage{amsmath,amssymb,amsthm,mathtools}
\usepackage{geometry}
\usepackage{hyperref}
\usepackage{xcolor}


\geometry{margin=1in}
\hypersetup{colorlinks=true,linkcolor=blue!50!black,citecolor=blue!50!black,urlcolor=blue!50!black}
\allowdisplaybreaks

% https://tex.stackexchange.com/questions/319398/making-the-dot-in-the-theorem-normal-but-not-as-like-bf
\makeatletter
\def\thmhead@plain#1#2#3{%
  \thmname{#1}\thmnumber{\@ifnotempty{#1}{ }\@upn{#2}}%
  \thmnote{ \the\thm@notefont(#3)}}% no group inside \thmnote
\let\thmhead\thmhead@plain
\makeatother

\newtheorem{theorem}{Theorem}
\newtheorem{corollary}[theorem]{Corollary}
\newtheorem{proposition}[theorem]{Proposition}
\newtheorem{lemma}[theorem]{Lemma}

\theoremstyle{definition}
\newtheorem{definition}[theorem]{Definition}
\newtheorem{example}[theorem]{Example}

\theoremstyle{remark}
\newtheorem{remark}[theorem]{Remark}

\DeclareMathOperator{\diam}{diam}

\title[Erd\H{o}s's diameter conjecture for separated distances fails in high dimensions]%
{Erd\H{o}s's diameter conjecture for separated distances fails in high dimensions}
\author{Boon Suan Ho}
\date{}

\begin{document}


\begin{abstract}
Erd\H{o}s asked whether every $n$-point set in Euclidean space whose
$\binom{n}{2}$ pairwise distances are mutually at least $1$ apart must have
diameter at least $(1+o(1))n^2$. We disprove this statement by constructing 
for every prime power $q$ a set
$\mathcal X_q\subset \mathbb R^{q^2+q}$ of $n=q+1$ points such that all
pairwise distances in $\mathcal X_q$ are mutually at least $1$ apart, while
\[
\diam(\mathcal X_q)
\le \Bigl(1-\frac{1}{\pi^2}+o(1)\Bigr)n^2.
\]
The proof is fully formalized in Lean 4.
\end{abstract}
\maketitle
\section{Introduction}
In a chapter of unsolved problems, Erd\H{o}s asked for lower bounds on the
diameter of an $n$-point set in Euclidean space when all pairwise distances are
separated from one another by at least $1$ \cite[Problem 20]{Erdos1997} 
(see also \cite[Problem 670]{Bloom}), and conjectured that
\begin{equation}\label{eq:erdos-conjecture}
\diam(\mathcal C) \ge (1+o(1))n^2,
\end{equation}
where the bound is independent of the dimension. He also gave a proof
in dimension $1$.

To the best of our knowledge, the exact higher-dimensional form of
Erd\H{o}s's question has received relatively little direct attention. The
closest earlier work seems to be Brass's study of the planar
``Erd\H{o}s-diameter'' \cite{Brass1996}, which assumes that distinct positive
distances are separated by at least $1$ and asks for the minimum possible
diameter. There is also a substantial literature on nearly equal distances,
beginning with Erd\H{o}s--Makai--Pach--Spencer \cite{EMPS1991} and
Erd\H{o}s--Makai--Pach \cite{EMP1993}, in which many pairwise distances are
allowed to lie in one or several short intervals. See also 
\cite{FranklKupavskii2023, PRV2006}. These problems are different
from the present one, but they are motivated by the same general question:
how strongly does the spacing of the distance set constrain the geometry of the
underlying point configuration?

The purpose of this note is to show that the conjectured bound
\eqref{eq:erdos-conjecture} is false in general:

\begin{theorem}\label{thm:main}
Let
\[
 c_*\coloneqq1-\frac{1}{\pi^2}=0.89867\dots.
\]
For infinitely many $n$, there exists an $n$-point set
$\mathcal X_n\subset \mathbb R^{n^2-n}$ such that
\begin{equation}\label{eq:sep-condition}
\bigl|\,|x-y|-|x'-y'|\,\bigr|\ge 1
\end{equation}
for every two distinct pairs
$\{x,y\}\neq \{x',y'\}$ from $\mathcal X_n$, while
\[
\diam(\mathcal X_n)\le (c_*+o(1))n^2.
\]
\end{theorem}

\newpage
The construction is explicit. Let $q$ be a prime power and put
\[
 m=q^2+q+1,
 \qquad n=q+1,
 \qquad N=\frac{m-1}{2}=\binom{n}{2}.
\]
Singer's theorem gives a cyclic $(m,n,1)$-difference set
$D\subset \mathbb Z/m\mathbb Z$ \cite{Singer1938}. Thus the $N$ unordered pairs from $D$
are indexed by the cyclic separations $1,\dots,N$. We then realize a chosen
distance profile $d_1<\dots<d_N$ in $\mathbb R^{2N}$ by placing the points on a
weighted product of regular $m$-gons. The profile is chosen so that the gaps
$d_{s+1}-d_s$ decrease with $s$; then the smallest gap is the last one, and
scaling by its reciprocal makes all pairwise distances at least $1$ apart.

\section{Singer difference sets and cyclic profiles}

We begin with the classical cyclic difference set construction.

\begin{proposition}[Singer \cite{Singer1938}]\label{prop:singer}
Let $q$ be a prime power and let $m=q^2+q+1$. There exists a set
$D\subset \mathbb Z/m\mathbb Z$ with $|D|=q+1$ such that every nonzero residue class in
$\mathbb Z/m\mathbb Z$ has a unique representation in the form $a-b$ with $a,b\in D$.
\end{proposition}

Fix such a set $D$, and write $n=q+1$ and $N=(m-1)/2$. For an unordered pair $\{t,u\}\subset \mathbb Z/m\mathbb Z$ with $t\neq u$, its
\emph{cyclic separation} is the unique integer $s\in[N]\coloneqq\{1,\dots,N\}$ such that
$t-u\equiv \pm s \pmod m$.

\begin{lemma}\label{lem:separations}
For each $s\in[N]$, there is exactly one unordered pair
$\{a,b\}\subset D$ such that
\[
a-b\equiv \pm s \pmod m.
\]
Consequently, the $\binom n2=N$ unordered pairs from $D$ are indexed by
$s=1,\dots,N$.
\end{lemma}

\begin{proof}
By Proposition \ref{prop:singer}, for each nonzero residue $r\in\mathbb Z/m\mathbb Z$,
there is a unique ordered pair $(a,b)\in D\times D$ with $a-b\equiv r\pmod m$.
The residues $s$ and $-s$ correspond to the same unordered pair, with the order
reversed. Thus each $s\in[N]$ determines at most one unordered pair.
Since there are $N$ values of $s$ and exactly $\binom n2=N$ unordered pairs,
each value occurs exactly once.
\end{proof}

\begin{proposition}\label{prop:profile-realization}
Let $W_1,\dots,W_N\ge 0$. For $t\in\mathbb Z/m\mathbb Z$ define
\[
 v_t\coloneqq\Bigl(\sqrt{W_r/2}\cos(2\pi rt/m),\ \sqrt{W_r/2}\sin(2\pi rt/m)\Bigr)_{1\le r\le N}
 \in \mathbb R^{2N}.
\]
If $\{t,u\}$ has cyclic separation $s\in[N]$, then
\begin{equation}\label{eq:distance-formula-general}
\|v_t-v_u\|^2
=\sum_{r=1}^N W_r\Bigl(1-\cos\frac{2\pi rs}{m}\Bigr).
\end{equation}
Hence, the multiset of pairwise distances among $\{v_t:t\in D\}$ is exactly the
collection of values
\[
 d_s\coloneqq\left(\sum_{r=1}^N W_r\Bigl(1-\cos\frac{2\pi rs}{m}\Bigr)\right)^{1/2},
 \qquad s=1,\dots,N,
\]
with each value (not necessarily distinct) occurring once.
\end{proposition}

\begin{proof}
In the $r$th two-dimensional factor, the squared distance between the $t$th and
$u$th points equals
\begin{align*}
&\frac{W_r}{2}\left(\bigl(\cos(2\pi rt/m)-\cos(2\pi ru/m)\bigr)^2
+\bigl(\sin(2\pi rt/m)-\sin(2\pi ru/m)\bigr)^2\right) \\
&\qquad= W_r\Bigl(1-\cos\frac{2\pi r(t-u)}{m}\Bigr).
\end{align*}
Summing over $r$ and using
the identity $\cos(a-b)=\cos(a)\cos(b)+\sin(a)\sin(b)$ yields \eqref{eq:distance-formula-general}. The final statement
follows from Lemma \ref{lem:separations}.
\end{proof}

\section{A one-frequency perturbation}

Fix $\varepsilon\coloneqq1/8$ and define
\[
 c_1\coloneqq1-\varepsilon=\frac78
 \quad\text{and}\quad
 c_j\coloneqq\frac1{j^2}
\]
for odd $j\ge3$. For $r\in[N]$, set
\begin{equation}\label{eq:Wr}
W_r\coloneqq\sum_{\substack{\text{odd }j\ge 1\\ j\equiv \pm r\!\!\!\pmod m}} c_j.
\end{equation}
Then $W_r>0$ for every $r$ (indeed, because $m$ is odd, exactly one of $r$ and
$m-r$ is an odd positive integer congruent to $\pm r \pmod m$, so the sum in
\eqref{eq:Wr} contains at least one positive term). Applying Proposition
\ref{prop:profile-realization} with these weights gives a point set
\[
\mathcal Y_m\coloneqq\{v_t:t\in D\}\subset \mathbb R^{2N}=\mathbb R^{m-1}.
\]

\begin{proposition}\label{prop:discrete-profile}
Let
\begin{equation}\label{eq:H-def}
H(\theta)\coloneqq\sum_{\text{odd }j\ge 1} c_j\,(1-\cos(j\theta)),
\qquad 0\le \theta\le \pi,
\end{equation}
and write
\[
G(\theta)\coloneqq\sqrt{H(\theta)}\qquad (0<\theta\le \pi).
\]
Then the pairwise distances in $\mathcal Y_m$ are exactly
\begin{equation}\label{eq:ds-grid}
 d_s=G(2\pi s/m), \qquad s=1,\dots,N.
\end{equation}
Moreover,
\begin{equation}\label{eq:H-explicit}
H(\theta)=\frac{\pi\theta}{4}-\varepsilon(1-\cos\theta)
\qquad (0\le \theta\le \pi).
\end{equation}
\end{proposition}

\begin{proof}
If $j\equiv \pm r\pmod m$, then for every integer $s$ we have
\[
\cos\frac{2\pi js}{m}=\cos\frac{2\pi rs}{m}.
\]
Therefore, for each $s\in[N]$,
\[
\sum_{r=1}^N W_r\Bigl(1-\cos\frac{2\pi rs}{m}\Bigr)
=\sum_{\substack{\text{odd }j\ge1\\ j\not\equiv 0\, (\mathrm{mod}\,m)}} c_j\Bigl(1-\cos\frac{2\pi js}{m}\Bigr).
\]
The terms with $j\equiv 0\pmod m$ contribute $0$, so by absolute convergence
this equals
\[
\sum_{\substack{j\ge 1\\ j\text{ odd}}} c_j\Bigl(1-\cos\frac{2\pi js}{m}\Bigr)
=H(2\pi s/m),
\]
which proves \eqref{eq:ds-grid}.
To obtain \eqref{eq:H-explicit}, write
\[
F(\theta)\coloneqq\sum_{k\ge 1}\frac{1-\cos(k\theta)}{k^2}
=\frac{\pi\theta}{2}-\frac{\theta^2}{4}
\qquad (0\le \theta\le 2\pi)
\]
using a standard Fourier series evaluation. The odd part is
\[
\sum_{\text{odd }j\ge 1}\frac{1-\cos(j\theta)}{j^2}
=F(\theta)-\frac14F(2\theta)
=\frac{\pi\theta}{4}
\qquad (0\le \theta\le \pi).
\]
Since the coefficients $c_j$ differ from $1/j^2$ only at $j=1$, equation
\eqref{eq:H-explicit} follows.
\end{proof}

\begin{lemma}\label{lem:monotone-concave}
The function $G$ is strictly increasing and strictly concave on $(0,\pi)$.
\end{lemma}

\begin{proof}
Set $a\coloneqq\pi/4$, so that
\[
H(\theta)=a\theta-\varepsilon(1-\cos\theta).
\]
We have
\[
H'(\theta)=a-\varepsilon\sin\theta\ge a-\varepsilon=\frac{\pi}{4}-\frac18>0,
\]
so $H$, and hence $G$, is strictly increasing on $(0,\pi)$.

For concavity we use
\begin{equation}\label{eq:G-second}
G''(\theta)=\frac{2H(\theta)H''(\theta)-H'(\theta)^2}{4H(\theta)^{3/2}}.
\end{equation}
Since $H''(\theta)=-\varepsilon\cos\theta$, a short calculation gives
\begin{equation}\label{eq:numerator-factorized}
2H(\theta)H''(\theta)-H'(\theta)^2
=-a^2+2a\varepsilon\bigl(\sin\theta-\theta\cos\theta\bigr)
-\varepsilon^2(1-\cos\theta)^2.
\end{equation}
Now set
\[
\phi(\theta)\coloneqq\sin\theta-\theta\cos\theta.
\]
Then
\[
\phi'(\theta)=\theta\sin\theta>0
\qquad (0<\theta<\pi),
\]
so $\phi$ is strictly increasing on $(0,\pi)$ and therefore
\[
0<\phi(\theta)<\phi(\pi)=\pi.
\]
Because $\varepsilon=1/8=a/(2\pi)$, equation
\eqref{eq:numerator-factorized} implies
\[
2H(\theta)H''(\theta)-H'(\theta)^2
\le -a^2+2a\varepsilon\pi-\varepsilon^2(1-\cos\theta)^2
=-\varepsilon^2(1-\cos\theta)^2<0
\]
for every $0<\theta<\pi$. By \eqref{eq:G-second}, this proves that $G$ is
strictly concave.
\end{proof}

\begin{corollary}\label{cor:smallest-gap}
The sequence
\[
 d_s=G(2\pi s/m),\qquad s=1,\dots,N,
\]
is strictly increasing and satisfies
\[
 d_{s+1}-d_s > d_{s+2}-d_{s+1}
 \qquad (1\le s\le N-2).
\]
In particular, the smallest gap is $d_N-d_{N-1}$.
\end{corollary}

\begin{proof}
By Lemma \ref{lem:monotone-concave}, $d_s$ is strictly increasing.
Since $G$ is strictly concave, its forward differences along any arithmetic 
progression are strictly decreasing, so $d_{s+1}-d_s>d_{s+2}-d_{s+1}$.
\end{proof}

Let
\[
\lambda_m\coloneqq\frac{1}{d_N-d_{N-1}},
\qquad
\mathcal X_m\coloneqq\lambda_m\mathcal Y_m.
\]
By Corollary \ref{cor:smallest-gap}, every two distinct distances in
$\mathcal X_m$ differ by at least $1$. Thus $\mathcal X_m$ satisfies
\eqref{eq:sep-condition}, and
\begin{equation}\label{eq:diam-lambda}
\diam(\mathcal X_m)=\lambda_m d_N=\frac{d_N}{d_N-d_{N-1}}.
\end{equation}

\section{Asymptotic diameter}

We now evaluate \eqref{eq:diam-lambda}.

\begin{proposition}\label{prop:asymptotic-diameter}
Let $m=q^2+q+1$ and $n=q+1$, with $q$ a prime power. Then
\[
\diam(\mathcal X_m)=\left(1-\frac{1}{\pi^2}+o(1)\right)m
=\left(1-\frac{1}{\pi^2}+o(1)\right)n^2.
\]
\end{proposition}

\begin{proof}
Set
\[
\theta_N\coloneqq\frac{2\pi N}{m}=\pi-\frac{\pi}{m},
\qquad
\theta_{N-1}\coloneqq\frac{2\pi(N-1)}{m}=\pi-\frac{3\pi}{m}.
\]
By the mean value theorem,
\[
 d_N-d_{N-1}=G(\theta_N)-G(\theta_{N-1})
 =\frac{2\pi}{m}G'(\xi_m)
\]
for some $\xi_m\in(\theta_{N-1},\theta_N)$. Since $\xi_m\to\pi$ and
$\theta_N\to\pi$, equation \eqref{eq:diam-lambda} yields
\[
\frac{\diam(\mathcal X_m)}{m}
=\frac{G(\theta_N)}{2\pi G'(\xi_m)}
\longrightarrow \frac{G(\pi)}{2\pi G'(\pi^-)},
\]
where $G'(\pi^-)\coloneqq\lim_{x\to\pi^-}G'(x)$. Now
\[
H(\pi)=\frac{\pi^2}{4}-2\varepsilon=\frac{\pi^2-1}{4},
\qquad
H'(\pi)=\frac{\pi}{4},
\]
so
\[
G'(\pi^-)=\frac{H'(\pi)}{2G(\pi)}=\frac{\pi}{8G(\pi)}.
\]
Therefore
\[
\frac{G(\pi)}{2\pi G'(\pi^-)}
=\frac{4H(\pi)}{\pi^2}
=1-\frac{1}{\pi^2}.
\]
Since $m=n^2-n+1$, the same asymptotic constant holds with $m$ replaced by
$n^2$.
\end{proof}

\begin{proof}[Proof of Theorem \ref{thm:main}]
Take $n=q+1$ with $q$ a prime power and construct $\mathcal X_m$ as above. Then
$\mathcal X_m\subset \mathbb R^{m-1}=\mathbb R^{n^2-n}$ has $n$ points,
satisfies \eqref{eq:sep-condition}, and has diameter bounded as in Proposition%
~\ref{prop:asymptotic-diameter}. Since there are infinitely many prime powers,
this provides infinitely many such values of $n$. The strict inequality
$c_*<1$ contradicts the dimension-free lower bound \eqref{eq:erdos-conjecture}.
\end{proof}

\section{Remarks}

\begin{remark}
Our construction shows that Erd\H{o}s's conjectured lower bound cannot hold
uniformly in the ambient dimension. It remains open whether, for each fixed
\(d\ge2\), every \(n\)-point set in \(\mathbb R^d\) whose pairwise distances
are mutually at least \(1\) apart must satisfy
\[
\diam(\mathcal C)\ge (1+o_d(1))n^2.
\]
\end{remark}

\begin{remark}
More generally, replacing $c_1=7/8$ by $1-\varepsilon$ and keeping
$c_j=1/j^2$ for odd $j\ge3$, the same construction yields
\[
H(\theta)=\frac{\pi\theta}{4}-\varepsilon(1-\cos\theta)
\]
and asymptotic constant $1-8\varepsilon/\pi^2$.
The choice $\varepsilon=1/8$ is a convenient round value at the edge of the
easy concavity argument.
In fact, within this one-frequency family, concavity holds for all
\[
0<\varepsilon\le \varepsilon_*:=\frac{\pi}{16}\bigl(\pi-\sqrt{\pi^2-4}\bigr),
\]
which gives the better constant
\[
1-\frac{8\varepsilon_*}{\pi^2}
=\frac{\pi+\sqrt{\pi^2-4}}{2\pi}
=0.885589\ldots.
\]
We have not attempted to optimize the coefficients $c_j$.
Numerically, optimizing several low odd frequencies suggests that the constant
for this particular construction can be lowered further, to about $0.85411$.
\end{remark}


\section*{Acknowledgements}
GPT-5.4 Pro was used to discover the construction of this paper, and Harmonic Aristotle was used to formalize the proof in Lean 4, with some assistance from GPT-5.4 Pro. 
All mathematical arguments and claims in the final manuscript were independently
verified by the author, who takes full responsibility for the paper.  
The formalization is available from \url{https://github.com/boonsuan/erdos670}.

The author thanks Way Yan Win and Alyxia Seah for helpful comments on a draft of this paper.

\begin{thebibliography}{99}

\bibitem{Bloom}
Thomas F. Bloom, \emph{Erd\H{o}s Problems}, \url{https://www.erdosproblems.com}.

\bibitem{Brass1996}
Peter Brass,
On the Erd\H{o}s-diameter of sets,
\emph{Discrete Mathematics} \textbf{150} (1996), 415--419.

\bibitem{Erdos1997}
Paul Erd\H{o}s,
Some unsolved problems,
in \emph{Combinatorics, Geometry and Probability: A Tribute to Paul
Erd\H{o}s}, edited by B.~Bollob\'as and A.~Thomason,
Cambridge University Press, Cambridge, 1997, 1--10.

\bibitem{EMP1993}
Paul Erd\H{o}s, Endre Makai, J\'anos Pach,
Nearly equal distances in the plane,
\emph{Combinatorics, Probability and Computing} \textbf{2} (1993), 401--408.

\bibitem{EMPS1991}
Paul Erd\H{o}s, Endre Makai, J\'anos Pach, Joel Spencer,
Gaps in difference sets, and the graph of nearly equal distances,
in \emph{Applied Geometry and Discrete Mathematics: The Victor Klee Festschrift},
edited by P.~Gritzmann and B.~Sturmfels,
DIMACS Series in Discrete Mathematics and Theoretical Computer Science
\textbf{4}, American Mathematical Society, Providence, RI, 1991, 265--273.

\bibitem{FranklKupavskii2023}
N\'ora Frankl, Andrey Kupavskii,
Nearly $k$-Distance Sets,
\emph{Discrete \& Computational Geometry} \textbf{70} (2023), 455--494.

\bibitem{PRV2006}
J\'anos Pach, Rado\v{s} Radoi\v{c}i\'c, Jan Vondr\'ak,
On the diameter of separated point sets with many nearly equal distances,
\emph{European Journal of Combinatorics} \textbf{27} (2006), 1321--1332.

\bibitem{Singer1938}
James Singer,
A theorem in finite projective geometry and some applications to
number theory,
\emph{Transactions of the American Mathematical Society} \textbf{43}
(1938), 377--385.

\end{thebibliography}

\end{document}
